\section{Complete-shell blow-up and observable removal}
\label{app:whole-shell}

We expand the proof of \cref{thm:whole-shell}.  The main issue is uniformity
over the full energy surface.  Local convergence near the harmonic orbit
would not rule out a second family that approaches the bottom along a
different configuration direction.

\subsection{Uniform compactness}

On \(h_a=2\pi+\epsilon^2\), positivity of the kinetic and potential excesses
gives
\begin{equation}
 |p|\le\sqrt2\epsilon,
 \qquad
 |\Psi_a(q)|^2
 \le\frac{1}{\pi}\log\left(1+\frac{\epsilon^2}{2\pi}\right).
 \label{eq:whole-shell-basic-bound}
\end{equation}
Write \((u,v)=\Psi_a(q)\).  The inverse formula yields
\begin{equation}
 q=(v,-c_av-av^2-u).
 \label{eq:inverse-bound-use}
\end{equation}
Both \(u\) and \(v\) are \(O(\epsilon)\), hence \(q=O(\epsilon)\)
uniformly over the entire shell.  Therefore the scaled states \((Q,P)\) lie
in a fixed compact set.  The removable extension of
\cref{eq:normalized-k} is smooth on a neighborhood of that set, and
\begin{equation}
 K_\epsilon\to K_0\quad\text{in }C^2,
 \qquad
 X_{K_\epsilon}\to X_{K_0}\quad\text{in }C^1.
 \label{eq:flow-convergence}
\end{equation}
Bounded-time flows and their first derivatives consequently converge
uniformly.

\subsection{Limiting returns}

Suppose \(Z_n\in\{K_{\epsilon_n}=1\}\), \(\epsilon_n\downarrow0\), and
\(\Phi_{\epsilon_n}^{T_n}(Z_n)=Z_n\) with \(0<T_n\le0.75\).  Compactness
gives \(Z_n\to Z_0\in\{K_0=1\}\) and \(T_n\to T_*\in[0,0.75]\) after a
subsequence.  If \(T_*=0\), then
\begin{equation}
 0=\frac{\Phi_{\epsilon_n}^{T_n}(Z_n)-Z_n}{T_n}
 =\frac{1}{T_n}\int_0^{T_n}
 X_{K_{\epsilon_n}}(\Phi_{\epsilon_n}^s(Z_n))\dd s
 \longrightarrow X_{K_0}(Z_0).
 \label{eq:no-small-return}
\end{equation}
The harmonic vector field vanishes only at the origin, which is not on
\(K_0=1\); therefore \(T_*>0\).

The slow harmonic component has first return \(T_-^0>0.75\), and the fast
component has returns at the multiples of \(T_+^0\), with
\(T_+^0<0.75<2T_+^0\).  Hence
\begin{equation}
 T_*=T_+^0,
 \qquad
 Z_0\in\Gamma_0:=\{K_0=1,Q_-=P_-=0\}.
 \label{eq:limiting-fast-circle}
\end{equation}
The set \(\Gamma_0\) is one phase circle, not a family of geometrically
different fast orbits.

\subsection{A phase-fixed local return}

At the positive fast turning point
\begin{equation}
 z_*=(Q_-=P_-=P_+=0, Q_+=\sqrt2/\omega_+),
 \label{eq:fast-turning-point}
\end{equation}
one has
\begin{equation}
 \partial_{Q_+}K_0(z_*)=\sqrt2\omega_+\ne0,
 \qquad
 \dot P_+(z_*)=-\sqrt2\omega_+\ne0.
 \label{eq:turning-transversality}
\end{equation}
The first inequality solves the energy equation for \(Q_+\), and the second
makes \(P_+=0\) a transverse event.  The varying sections can therefore be
identified by \((Q_-,P_-)\).  In normalized slow coordinates
\begin{equation}
 x_-=\sqrt{\omega_-}Q_-,
 \qquad y_-=P_-/\sqrt{\omega_-},
 \label{eq:normalized-slow}
\end{equation}
the limiting derivative is the rotation
\begin{equation}
 D\Pi_0(z_*)=R_{2\pi/\rho_a}.
 \label{eq:limiting-return-rotation}
\end{equation}
Its fixed-point derivative is invertible because
\(4\sin^2(\pi/\rho_a)>0\).  The parameter-dependent implicit-function
theorem gives one nearby fixed point for every sufficiently small
\(\epsilon\).

\begin{lemma}[Unique oriented local crossing]
\label{lem:unique-local-crossing}
There are a section neighborhood \(U\) of \(z_*\), a number \(\eta>0\),
and \(\epsilon_0>0\) with the following property.  Let
\(0<\epsilon_n<\epsilon_0\), and let primitive periodic trajectories
\(z_n(t)\) have periods \(T_n\to T_+^0\).  If, after phase alignment,
\(z_n\to z_0\) in \(C^1\) over one period, where \(z_0\) parametrizes
\(\Gamma_0\) and \(z_0(0)=z_*\), then \(z_n\) has exactly one intersection
per primitive circuit with
\[
 U\cap\{K_{\epsilon_n}=1,\ P_+=0,\ Q_+>0\},
\]
and the event has the same orientation as at \(z_*\).
\end{lemma}

\begin{proof}
Parametrize the limiting circle on the cyclic interval
\([ -\eta,T_+^0-\eta]\).  Since \(P_+\) has a simple zero at the positive
turning point, choose \(\eta\), \(c>0\), and \(U\) so that
\(\dot P_+\le-c\) on the part of \(\Gamma_0\) in \(U\), the values of
\(P_+\) at times \(-\eta\) and \(\eta\) have opposite signs, and the
remaining compact arc of \(\Gamma_0\) is disjoint from the closure of the
local section in \(U\).  Bounded-time \(C^1\) convergence preserves the two
endpoint signs and gives \(\dot P_+\le-c/2\) whenever \(z_n(t)\in U\).
The intermediate-value theorem supplies one local zero, strict monotonicity
supplies at most one, and uniform position convergence excludes another
intersection on the complementary arc.  Identifying the endpoints of the
period interval counts the crossing only once.
\end{proof}

\subsection{Primitivity and globalization}

Suppose a sequence of the returns above is an \(m_n\)-fold iterate of a
primitive period \(\tau_n\).  If \(m_n\to\infty\), then
\(\tau_n\to0\), contradicting \cref{eq:no-small-return}.  Along a subsequence
with \(m_n=m\), flow convergence would give
\begin{equation}
 \Phi_0^{T_+^0/m}(Z_0)=Z_0.
 \label{eq:fractional-return}
\end{equation}
On \(\Gamma_0\), this is possible only for \(m=1\).  The return is therefore
primitive.

If a second geometric orbit existed for every member of a sequence
\(\epsilon_n\downarrow0\), the limiting classification and first-derivative
flow convergence would make its phase-aligned trajectory converge uniformly
in \(C^1\) to \(\Gamma_0\) over one period.  By
\cref{lem:unique-local-crossing}, it has exactly one oriented crossing of the
local section neighborhood per primitive circuit.  Its return after that
circuit is therefore the next local return, so the crossing is a fixed point
of \(\Pi_{\epsilon_n}\).  The implicit-function theorem provides only the
continued fast fixed point.  The contradiction proves complete-shell
uniqueness.

The globalization step is exactly what permits the symbol in the trace
formula to be \(A=1\).  Before this step, an observable supported near the
fast orbit would isolate a legitimate microlocal term, but the spectral sum
would contain matrix elements
\(\langle\psi_k,A_\h\psi_k\rangle\).  After the full return classification,
the observable is unnecessary and the result becomes eigenvalue-only.
